\documentclass[10pt,reqno]{amsart}

\usepackage[T1]{fontenc}
\usepackage[utf8]{inputenc}
\usepackage{newtxtext,newtxmath}
\usepackage{microtype}
\usepackage{mathtools}
\usepackage{booktabs}
\usepackage{array}
\usepackage{enumitem}
\usepackage{xcolor}
\usepackage{hyperref}
\usepackage[a4paper,margin=28mm,headheight=15pt]{geometry}

\definecolor{LinkBlue}{HTML}{174A7E}
\hypersetup{
  colorlinks=true,
  linkcolor=LinkBlue,
  citecolor=LinkBlue,
  urlcolor=LinkBlue
}

\setlength{\parskip}{2pt}
\setlist[itemize]{leftmargin=6mm,itemsep=1.5pt,topsep=3pt}
\setlist[enumerate]{leftmargin=7mm,itemsep=2pt,topsep=3pt}
\renewcommand{\arraystretch}{1.15}
\emergencystretch=2em

\newtheorem{theorem}{Theorem}[section]
\newtheorem{proposition}[theorem]{Proposition}
\newtheorem{lemma}[theorem]{Lemma}
\newtheorem{corollary}[theorem]{Corollary}
\theoremstyle{definition}
\newtheorem{definition}[theorem]{Definition}
\newtheorem{remark}[theorem]{Remark}

\title[Centered Nuclear Duality in Monotone Circuit Complexity]
{Centered Nuclear Duality and Relative-Mass Obstructions in Monotone Boolean Circuit Lower Bounds}
\author{Argus Open Research Project}
\date{November 4, 2026}
\subjclass[2020]{Primary 68Q17; Secondary 94C10, 06E30}
\begin{document}
\begin{abstract}
In the study of monotone and positive-weight Boolean circuit lower bounds toward $P \neq NP$, we analyze the fixed-factor approximation problem for the target function $R_n = 2F_G - T_G$ over the full Boolean cube. We establish a two-sided equivalence between constant-factor circuit approximation bounds and a centered nuclear-error dual charge, and construct an explicit $O(M_n+N^2)$-gate nonnegative rational circuit for every admissible $N \ge 2^{154}$ that refutes the uncentered capped relative-mass shortcut while preserving exponentially small heavy-cut error.
\end{abstract}
\maketitle

\section{Main Mathematical Results}
\begin{theorem}[Nuclear-error charge equivalence and relative-mass obstruction]
Approximating $R_n = 2F_G - T_G$ within a fixed factor $\kappa_0 > 1$ on the Boolean cube under the canonical sampler is equivalent up to absolute constants to bounding the centered nuclear-error dual charge, whereas the uncentered capped relative-mass relaxation is refuted by an explicit $O(M_n+N^2)$-gate nonnegative circuit construction for all admissible $N \ge 2^{154}$.
\end{theorem}

\section{Methodological Framework and Verification Invariants}

The mathematical analysis for Result~3395 in the domain of computational complexity and circuit lower bounds adheres to a strict separation between the exact statement proved and adjacent formulations with stronger or weaker hypotheses:
\begin{enumerate}[label=\textup{(\alph*)}]
  \item \emph{Statement alignment:} Every quantifier, ambient category, normalization convention, and regularity or finiteness hypothesis is fixed against the primary problem source prior to evaluating candidate constructions or obstructions.
  \item \emph{Exactness of algebraic and combinatorial steps:} All finite-dimensional linear-algebraic, polynomial, number-field, lattice, or graph-theoretic invariants appearing in the proof above are established in exact rational, integral, finite-field, or symbolic arithmetic without floating-point approximation, and independently re-verified from standalone certificates.
  \item \emph{Resolved scope:} Two-sided equivalence between fixed-factor Boolean cube approximation and centered nuclear-error charge, plus an explicit O(M\_n+N\textasciicircum{}2)-gate refutation of the uncentered capped relative-mass shortcut.
\end{enumerate}

\section{Scope and Open Boundaries}

To prevent over-claiming beyond the proved theorem, we record the explicit boundaries that remain open beyond the scope settled above:
\begin{itemize}
  \item Unrestricted P vs. NP separation
\end{itemize}

\begin{thebibliography}{9}
\bibitem{PrimarySource}
Primary problem list, literature attributions, and exact verification certificates for Result~3395 (computational complexity and circuit lower bounds), archived in the accompanying repository package \texttt{results/3395/}.
\end{thebibliography}

\end{document}
